My CR5 pipeline carried observations and commands through ROS2 collection, HDF5 synchronization, LeRobot v2.0 conversion, OpenPI configuration, WebSocket serving, and robot execution. I preserved seven-dimensional state and action semantics, yet each transition still depended on a separately implemented assumption. My next step is to specify and validate one end-to-end interface contract spanning recognition, learning, serving, and execution rather than treating those transitions independently. HKUST's MSc in Information Technology offers the software and learning depth for that work.

With the current courses available in my intake and CSIT6910 Independent Project approved, CSIT5100 Engineering Reliable Software Systems would guide the contract's reliability and testing logic. CSIT5410 Recognition Systems and CSIT5910 Machine Learning would strengthen how its perceptual inputs and learned outputs are defined, while CSIT6910 could provide the setting to implement the pipeline as one evaluated system. I would trace schemas, clocks, service responses, and execution timing through explicit interface and failure checks, so a discrepancy could be localized instead of absorbed into overall model behavior. Developing that contract-backed system would prepare me to engineer embodied-agent stacks whose behavior remains traceable from sensor representation to robot command in robotics R\&D.
